\chapter{Canonical Resolution and Diagnostic Invariance}\label{chap:3}

\section*{Overview of the Chapter}

The question of this chapter is \textbf{under what conditions the same obstruction can be
diagnosed when we merge information or repartition the regions that are read locally}.
In Chapter~\ref{chap:2} we built obstruction classes to gluing from local states and
transitions. To share that verdict across different resolutions, we must specify the
information to retain and the maps that compare local data.

For example, suppose that a state is a pair of a displayed value and an internal value,
and that the chosen Law reads only the displayed value. Merging states with the same
displayed value loses no Law value. If, however, we merge even the regions on which
local states are placed, the cycles of overlaps may change. Preserving Law values and
preserving obstructions on cycles each require conditions to be checked.

In the first half we construct the coarsest resolution that preserves the evaluation of
the Laws, and generate a diagnostic complex from its evaluation values. We then build an
actual cochain map from a comparison of covers, and prove a sufficient condition for it
to be an isomorphism on first cohomology. In the second half we reduce invariance over
all Law families for which both resolutions are adequate to a finite computation, and
connect the integer-coefficient obstructions of Chapter~\ref{chap:2} with the
rational-valued diagnoses of this chapter.

\begin{xltabular}{\linewidth}{@{}LLL@{}}
\toprule
Question & Construction & Outcome \\
\midrule
\endhead
How far can information be merged while the Laws remain readable? & Identify sources on which all Law values agree & Universality of the canonical resolution, and conditions for presentation by specified extraction methods (\S\S\ref{sec:3.1}--\ref{sec:3.2}) \\
\addlinespace[3pt]
How to compare changes of resolution and of local structure? & Build maps from coefficients per Law value and from comparisons of edges and faces & Maps carrying diagnostic classes, and isomorphism under Condition C (\S\S\ref{sec:3.3}--\ref{sec:3.5}) \\
\addlinespace[3pt]
Is the diagnosis preserved whichever Laws are chosen? & Compute kernels and cokernels for each subset of values & A necessary and sufficient condition for uniform invariance, finite decision, and a counterexample for local observation (\S\S\ref{sec:3.6}--\ref{sec:3.7}) \\
\addlinespace[3pt]
Does fixing the structure also fix the obstruction? & Examine the structural support and the chosen coefficients separately & Invariance of the structural nerve, and vanishing for per-source coefficients (\S\ref{sec:3.8}) \\
\addlinespace[3pt]
Does a vanishing diagnosis mean a vanishing actual obstruction? & Send integer generators to Law-value coordinates, and turn rational corrections into integer ones & Correspondence of specified obstruction classes and equivalence of vanishing (\S\S\ref{sec:3.9}--\ref{sec:3.10}) \\
\addlinespace[3pt]
Does refining the cover give the same verdict? & Use an actual refinement from three charts to four & The comparison square, invariance of the obstruction verdict, and concrete zero and nonzero examples (\S\ref{sec:3.11}) \\
\bottomrule
\end{xltabular}

\subsection*{Why compare diagnoses}

The same state of software may be read down to its internal values, or only through its
published values. Local consistency, too, can be examined component by component, or
over regions that group several components together. Comparing results while varying
such choices requires maps that preserve the meaning of diagnostic values.

The diagnosis of this chapter is a first cohomology built from Law values and overlap
information, together with specified classes in it. For the lenses and protocols of
Chapter~\ref{chap:1}, choosing the required readouts and results of operations as the
evaluations determines what the resolution retains. For the gluing problem of
Chapter~\ref{chap:2}, what matters is that the vanishing of a diagnostic class can be
brought back to the existence of local corrections.
\S\S\ref{sec:3.9}--\ref{sec:3.11} provide this way back, with the same inputs and
concrete maps.

\section{The canonical resolution preserving Law evaluations}\label{sec:3.1}

A resolution specifies which sources are distinguished. Here we represent by a
surjection the part of the choices carried by a reading of Chapter~\ref{chap:1} that
concerns this distinction. Local covers and coefficients are added from
\S\ref{sec:3.3} on.

\begin{definition}[Resolutions and Law evaluations]\label{def:3.1}

Let $S$ be the set of sources, and give a resolution as a surjection $q:S\to Q$.
Fix a finite Law index set $L$ and evaluations
$v_\ell:S\to\mathrm{Val}_\ell$.
These are the same kind of input as the Law evaluations used in Definition~\ref{def:2.35}.
When studying the equations of Chapter~\ref{chap:1}, chosen residuals and the like can
serve as evaluations.

We say that $q$ is $L$-adequate if for each $\ell\in L$ there exists
$\bar v_\ell:Q\to\mathrm{Val}_\ell$ with
\begin{equation*}
v_\ell=\bar v_\ell q
\tag{3.1}\label{eq:3.1}
\end{equation*}
This is the condition that ``the value of each Law can be read back from $q(s)$''.
Since the data involved differ from those of the adequacy of Chapter~\ref{chap:1} (that a cover
can read information and interactions), we call this notion $L$-adequacy to distinguish
the two.

For two resolutions we write
\begin{equation*}
q_c\preceq q_f
\quad\Longleftrightarrow\quad
\forall s,t\in S,\
q_f(s)=q_f(t)\Rightarrow q_c(s)=q_c(t)
\tag{3.2}\label{eq:3.2}
\end{equation*}
and say that $q_c$ is coarser than $q_f$.
Two sources identified on the finer side are identified on the coarser side as well.

\end{definition}

\begin{lemma}[Descent of evaluations and comparison maps]\label{lem:3.2}

The following hold.

\begin{enumerate}
\item A map $v:S\to V$ descends through $q$ if and only if
$q(s)=q(t)\Rightarrow v(s)=v(t)$. The descended map is unique.
\item If $q_c\preceq q_f$, there is a unique surjection $\pi:Q_f\to Q_c$ with
$q_c=\pi q_f$.
\item If both sides are $L$-adequate, then $\bar v_{\ell,c}\pi=\bar v_{\ell,f}$.
\end{enumerate}

\end{lemma}

\begin{proof}

The forward direction of 1 follows by composing maps. For the converse, take a
representative $s$ of $u\in Q$ and set $\bar v(u)=v(s)$. By the hypothesis this does
not depend on the representative, and by the surjectivity of $q$ it is unique.
For 2, apply 1 with $v=q_c$; the surjectivity of $q_c$ makes $\pi$ surjective as well.
For 3, both sides equal $v_\ell(s)$ at $q_f(s)$, so they agree by surjectivity.

\end{proof}

\begin{definition}[The canonical resolution]\label{def:3.3}

Define an equivalence relation on sources by
\begin{equation*}
s\sim_L t
\quad\Longleftrightarrow\quad
\forall\ell\in L,\ v_\ell(s)=v_\ell(t)
\tag{3.3}\label{eq:3.3}
\end{equation*}
We call the quotient map $q_L:S\to S/{\sim_L}$ the canonical resolution of this
Law family.

\end{definition}

\begin{theorem}[Universality of the canonical resolution]\label{thm:3.4}

The map $q_L$ is $L$-adequate, and for every $L$-adequate $q:S\to Q$ there is a
unique map $u_q:Q\to S/{\sim_L}$ satisfying
\begin{equation*}
u_q q=q_L
\tag{3.4}\label{eq:3.4}
\end{equation*}
Hence $q_L\preceq q$.
A resolution with this property is unique up to a unique bijection commuting with the
maps from the sources.

\end{theorem}

\begin{proof}

The assignment $[s]\mapsto v_\ell(s)$ does not depend on the representative by
\eqref{eq:3.3}, so $q_L$ is adequate. If $q$ is adequate, then $q(s)=q(t)$ implies
that all Law values agree, hence $s\sim_L t$. By Lemma~\ref{lem:3.2} the map $u_q$
exists and is unique. If two resolutions have the same universal property, we obtain
maps through each other. Both composites commute with the maps from the sources, and
by uniqueness they are the identities.

\end{proof}

The canonical resolution is the coarsest among the resolutions that can read the
chosen Laws. The distinctions to be retained are thus determined before any choice of
implementation or of local covers.

\begin{example}[Displayed and internal values]\label{ex:3.5}

Take $S=\{0,1\}^2$ and a single Law $v_\ell(v,h)=v$.
The canonical resolution has the same equivalence relation as $q_c(v,h)=v$, which
reads the first component. The identity $q_f(v,h)=(v,h)$ is a finer adequate
resolution, and the comparison map is $\pi(v,h)=v$.
The states $(v,0)$ and $(v,1)$ can be merged, while $v=0$ and $v=1$, which the Law
distinguishes, remain apart. Adding a Law that also reads the internal value makes the
canonical resolution agree with the identity.

\end{example}

\section{Finite computation and presentation by extraction methods}\label{sec:3.2}

That the quotient of the canonical resolution can be formed, and that an available
extraction method realizes this quotient, are verified separately. For finite sources,
both can be computed as comparisons of equivalence relations.

\begin{construction}[Computing the finite partition]\label{cons:3.6}

Suppose that $S$ is finite and that equality of evaluation values is decidable.
For each source, compute the finite tuple $(v_\ell(s))_{\ell\in L}$.
Let $\mathcal{P}_L$ be the family of nonempty subsets obtained by grouping sources
according to agreement of these tuples.
The map $s\mapsto\{t\in S\mid t\sim_L s\}$ is a surjection onto $\mathcal{P}_L$.
Equality of two outputs is equivalent to \eqref{eq:3.3}, so this computation presents
the canonical resolution. If the Law index set is empty and the set of sources is not,
everything merges into a single component.

\end{construction}

\begin{definition}[Presentation by permitted extraction methods]\label{def:3.7}

Let $\mathcal{D}$ be a specified family of extraction doctrines of
Chapter~\ref{chap:1} over a common source set $S$.
The resolution induced by $D\in\mathcal{D}$ is the quotient by
\begin{equation*}
s\sim_D t
\quad\Longleftrightarrow\quad
\mathrm{Atomize}_D(s)=\mathrm{Atomize}_D(t)
\tag{3.5}\label{eq:3.5}
\end{equation*}
When $\sim_D=\sim_L$ for some $D\in\mathcal{D}$, we say that the canonical resolution
is presentable by this family.
Given a finite $\mathcal{D}$ and a method for deciding equality of extracted families,
presentability can be decided by comparing \eqref{eq:3.3} and \eqref{eq:3.5} over all
pairs of sources.

\end{definition}

\begin{example}[Adequate extraction methods alone do not yield presentability]\label{ex:3.8}

Take six sources $a_0,a_1,a_2,b_0,b_1,b_2$, with Law value 0 on the $a_i$ and 1 on
the $b_i$. Each of the following rows is a partition of the sources on which the
extraction results agree.

\begin{xltabular}{\linewidth}{@{}LL@{}}
\toprule
Extraction method & Components of sources with the same result \\
\midrule
\endhead
$D_0$ & $\{a_0,a_1,a_2\},\ \{b_0,b_1,b_2\}$ \\
\addlinespace[3pt]
$D_a$ & $\{a_0\},\ \{a_1,a_2\},\ \{b_0,b_1,b_2\}$ \\
\addlinespace[3pt]
$D_b$ & $\{a_0,a_1,a_2\},\ \{b_0\},\ \{b_1,b_2\}$ \\
\bottomrule
\end{xltabular}

Assigning a distinct Atom to each component and taking its singleton as the
extraction result realizes these partitions by actual doctrines. In
Definition~\ref{def:1.3} it suffices to take the normalization to be the identity, to
define the semantics of a source as ``it is the assigned Atom'', and to take the
remaining permission predicates to be true.

All three preserve the Law values. In $\{D_0,D_a\}$, the doctrine $D_0$ presents the
canonical resolution. By contrast, each member of $\{D_a,D_b\}$ retains a distinction
that the Law does not need. The former separates $a_0,a_1$ and the latter
separates $b_0,b_1$, so a family permitting only these two cannot present the
canonical resolution.

\end{example}

\section{The diagnostic complex from Law values}\label{sec:3.3}

From now on, sources are finite and $q:S\to Q$ is $L$-adequate.
For each local region we take as coordinates the distinct Law values readable there.
The point of the construction is that coordinates are not duplicated as more sources
share the same value.

\begin{definition}[Finite nerves with supports]\label{def:3.9}

Let finite sets $N_0,N_1,N_2$ be the sets of charts, edges, and faces, respectively.
An edge $e$ has a start $s(e)$ and an end $t(e)$.
A face $f$ has, for some charts $\alpha,\beta,\gamma$, three edges
$e_0:\alpha\to\beta$, $e_1:\alpha\to\gamma$, and
$e_2:\beta\to\gamma$. Parallel edges, self-loops, and repeated vertices are allowed.

To each chart we assign a nonempty support $S_\alpha\subseteq Q$, and we form the
supports of edges and faces by
\begin{equation*}
S_e=S_{s(e)}\cap S_{t(e)},\qquad
S_f=S_{e_0}\cap S_{e_1}\cap S_{e_2}
\tag{3.6}\label{eq:3.6}
\end{equation*}
The support of an edge or a face may be empty.
This support represents the values of the resolution readable on each cell.
The existence of edges and faces is part of the chosen local comparison data.
When this is applied to a cover of a space $X$, an open set $U_\alpha\subseteq X$
represents the region on which local states are placed, and the support
$S_\alpha\subseteq Q$ represents the values of the resolution readable there.
Cells are taken from intersections of the open sets, and which Law values are read on
a cell is determined by its support.

\end{definition}

\begin{construction}[The Law-value diagnostic complex]\label{cons:3.10}

Let $\Lambda=\{(\ell,v_\ell(s))\mid\ell\in L,\ s\in S\}$ be the set of labels
occurring overall.
Define the labels occurring on a cell $\sigma$ by
\begin{equation*}
\Lambda_\sigma
=\{(\ell,\bar v_\ell(u))\mid \ell\in L,\ u\in S_\sigma\}
\subseteq\Lambda
\tag{3.7}\label{eq:3.7}
\end{equation*}
A label is a pair of a Law index and a value; the element $u$ realizing the value is
not part of it. This pair is the name of a coordinate, distinguished from the rational
number placed in that coordinate.
Define the diagnostic cochain groups by
\begin{equation*}
D_L^n(q,N)=\prod_{\sigma\in N_n}\mathbb{Q}^{\Lambda_\sigma}
\quad(n=0,1,2)
\tag{3.8}\label{eq:3.8}
\end{equation*}
and the differentials, between coordinates with the same label, by
\begin{equation*}
\begin{aligned}
(d^0b)_{e,\lambda}
&=b_{t(e),\lambda}-b_{s(e),\lambda},\\
(d^1z)_{f,\lambda}
&=z_{e_0,\lambda}-z_{e_1,\lambda}+z_{e_2,\lambda}
\end{aligned}
\tag{3.9}\label{eq:3.9}
\end{equation*}
A label occurring on an edge occurs on both endpoints, and a label occurring on a face
occurs on its three edges, so every term on the right-hand side is defined.

\end{construction}

\begin{lemma}[Differentials and the diagnostic group]\label{lem:3.11}

We have $d^1d^0=0$. Hence
\begin{equation*}
H^1_{\mathrm{diag}}(q,N;L)
=\ker d^1/\mathrm{im}\,d^0
\tag{3.10}\label{eq:3.10}
\end{equation*}
is defined.

\end{lemma}

\begin{proof}

Computing \eqref{eq:3.9} on one face and one label gives
$(b_\beta-b_\alpha)-(b_\gamma-b_\alpha)+(b_\gamma-b_\beta)=0$.

\end{proof}

Diagnostic classes arising from concrete local data are constructed as elements of
this group by means of the coefficient comparison of \S\ref{sec:3.9}.

\begin{proposition}[Decomposition by Law value]\label{prop:3.12}

For $\lambda\in\Lambda$, let $N_\lambda$ be the subnerve retaining exactly the cells
with $\lambda\in\Lambda_\sigma$. With its complex of constant rational coefficients
we have
\begin{equation*}
\begin{aligned}
D_L^\bullet(q,N)
&\cong\bigoplus_{\lambda\in\Lambda}C^\bullet(N_\lambda;\mathbb{Q}),\\
H^1_{\mathrm{diag}}(q,N;L)
&\cong\bigoplus_{\lambda\in\Lambda}H^1(N_\lambda;\mathbb{Q})
\end{aligned}
\tag{3.11}\label{eq:3.11}
\end{equation*}
\end{proposition}

\begin{proof}

Reorder the coordinates $(\sigma,\lambda)$ of \eqref{eq:3.8} with $\lambda$ first.
The sets are finite, so products and direct sums agree.
The differentials do not change labels, so the reordering commutes with them.
The cocycle condition and membership in the coboundaries also hold label by label,
giving an isomorphism on the quotients as well.

\end{proof}

By this decomposition, the comparison of diagnoses splits into comparisons of ``the
regions where a single Law value is visible''. For example, taking every chart support
in Example~\ref{ex:3.5} to be the whole set, one copy of the complex of the same nerve
appears for each of the two Law values 0 and 1.

\section{From comparisons of resolutions to cochain maps}\label{sec:3.4}

The map carrying the diagnosis of the coarser side to the finer side can be built
from the correspondence between evaluation values and local regions.
Note that the direction in which resolution values are sent is opposite to the
direction in which cochains are pulled back.
Below we also call the coarser side coarse and the finer side fine.

\begin{definition}[Local comparison data]\label{def:3.13}

Take $q_c\preceq q_f$ and the comparison map $\pi:Q_f\to Q_c$ of Lemma~\ref{lem:3.2}.
For the nerves with supports on the two sides, specify the following.

\begin{itemize}
\item A chart map $r:N_{f,0}\to N_{c,0}$.
\item A map $r_1$ sending each fine edge either to a coarse edge or to a symbol
$\bot$ denoting contraction.
\item A map $r_2$ sending each fine face either to a coarse face or to $\bot$.
\end{itemize}

A mapped edge preserves the start and the end, and a mapped face preserves the
numbering of its three edges.
The two endpoints of a contracted edge map to the same coarse chart.
For a contracted face, its three edges are all contracted as well.
We also assume
\begin{equation*}
\pi(S^f_\alpha)\subseteq S^c_{r(\alpha)}
\tag{3.12}\label{eq:3.12}
\end{equation*}
The corresponding inclusions of supports for edges and faces follow from
\eqref{eq:3.6} and the incidence relations.

\end{definition}

\begin{proposition}[The generated diagnostic comparison]\label{prop:3.14}

From the data of Definition~\ref{def:3.13} and $L$-adequacy on both sides, linear
maps $T^n:D_L^n(q_c,N_c)\to D_L^n(q_f,N_f)$ are given by
\begin{equation*}
\begin{aligned}
(T^0b)_{\alpha,\lambda}&=b_{r(\alpha),\lambda},\\
(T^1z)_{e,\lambda}
&=\begin{cases}
z_{r_1(e),\lambda}&r_1(e)\ne\bot,\\
0&r_1(e)=\bot,
\end{cases}\\
(T^2w)_{f,\lambda}
&=\begin{cases}
w_{r_2(f),\lambda}&r_2(f)\ne\bot,\\
0&r_2(f)=\bot
\end{cases}
\end{aligned}
\tag{3.13}\label{eq:3.13}
\end{equation*}
and satisfy $T^1d_c^0=d_f^0T^0$ and $T^2d_c^1=d_f^1T^1$.
The same maps therefore induce a linear map
\begin{equation*}
T_{\mathrm{diag}}:
H^1_{\mathrm{diag}}(q_c,N_c;L)
\longrightarrow H^1_{\mathrm{diag}}(q_f,N_f;L),
\qquad [z]\longmapsto[T^1z]
\tag{3.14}\label{eq:3.14}
\end{equation*}
Under the decomposition \eqref{eq:3.11}, this is the direct sum of the comparisons of
the individual $N_\lambda$.

\end{proposition}

\begin{proof}

Take $u$ realizing $\lambda=(\ell,v)$ on a fine cell.
On a mapped cell, \eqref{eq:3.12} and Lemma~\ref{lem:3.2} show that $\pi(u)$ realizes
the same label on the coarse side.
The coordinates in \eqref{eq:3.13} therefore exist.

On a mapped edge, $T^1d_c^0=d_f^0T^0$ follows from the condition that the endpoints
are preserved. On a contracted edge, the two chart values agree and their difference
is zero. On a mapped face, the numbering of the three edges is preserved, so
$T^2d_c^1=d_f^1T^1$. On a contracted face, all three edges are contracted and all
three terms are zero. The maps preserve cocycles and coboundaries, hence descend to
the quotients, and since they do not change labels they also commute with the
decomposition.

\end{proof}

The same argument applies to cochain maps with coefficients in abelian groups.
A comparison commuting with the differentials carries classes and sends zero classes
to zero classes. If, moreover, the comparison in each degree is an isomorphism, its
inverse also commutes with the differentials, and the comparison of classes is an
isomorphism. The standard \v{C}ech comparison by refinement of covers and restriction
of sections is given in
\cite[\href{https://stacks.math.columbia.edu/tag/09UY}{Tag 09UY}]{Stacks}.
In \S\ref{sec:3.11} we compute the restriction maps for the actual covers of
Chapter~\ref{chap:2}.

\section{Diagnostic invariance under Condition C}\label{sec:3.5}

Once the comparison map is built, the question is whether it loses no class and
captures every class on the fine side. We prove these two properties as injectivity
and surjectivity, respectively. To this end we examine the structure on the
fine side that gathers inside a single coarse chart.

\begin{definition}[Coordinate fibers and Condition C]\label{def:3.15}

Fix a label $\lambda$ and work in the subnerve of \eqref{eq:3.11}.
The fiber of a coarse chart $v$ consists of the fine charts with $r(\alpha)=v$,
together with all fine edges whose two endpoints lie in this fiber.
A face whose three edges all lie within the fiber is called an internal face.
The edges within a fiber include, besides contracted edges, edges mapping to coarse
self-loops.

Condition C consists of the following seven clauses.

\begin{xltabular}{\linewidth}{@{}LL@{}}
\toprule
Clause & Content \\
\midrule
\endhead
C0: agreement of supports & For each coarse chart $v$, $S^c_v=\bigcup_{r(\alpha)=v}\pi(S^f_\alpha)$ \\
\addlinespace[3pt]
C1: connectedness of fibers & For each chart of the coarse subnerve of each $\lambda$, the fiber is nonempty and connected as a graph with edge orientations forgotten \\
\addlinespace[3pt]
C2: lifting of edges & Each coarse edge and each of its labels admits a fine edge with the same label mapping to that edge \\
\addlinespace[3pt]
C3: cycles within fibers & Every rational 1-cycle within a fiber is a rational linear combination of boundaries of internal faces \\
\addlinespace[3pt]
C4: lifting of faces & Each coarse face and each of its labels admits a fine face with the same label mapping to that face \\
\addlinespace[3pt]
C5: uniqueness of edge lifts & At most one fine edge maps to any given coarse edge \\
\addlinespace[3pt]
C6: reflection of self-loops & A fine edge mapping to a coarse self-loop is itself a self-loop \\
\bottomrule
\end{xltabular}

C0, C5, and C6 are imposed on the whole nerve, and C1--C4 on every
$\lambda\in\Lambda$. The boundaries of chains in C3 are $\partial e=t(e)-s(e)$ and
$\partial f=e_0-e_1+e_2$. A 1-cycle is a finite sum of edges with rational
coefficients in which, at each vertex, the inflowing and the outflowing coefficients
agree.

\end{definition}

\begin{lemma}[Correction within fibers]\label{lem:3.16}

Assume C1 and C3 in the subnerve of a single label. For every 1-cocycle $z$ on the
fine side there is a 0-cochain $p$ such that $z-d_f^0p$ vanishes on all edges within
fibers.

\end{lemma}

\begin{proof}

Choose a root vertex within one fiber.
Let $p(\alpha)$ be the signed sum of $z$ along a path from the root to $\alpha$.
A path exists by C1. The difference of two paths is a 1-cycle, and by C3 a sum of
boundaries of internal faces.
A cocycle vanishes on the boundary of each face, so $p(\alpha)$ does not depend on
the choice of path.
For an edge $e:\alpha\to\beta$, appending $e$ to a path from the root to $\alpha$
gives $p(\beta)-p(\alpha)=z(e)$.
Carrying out this construction in each fiber defines $p$ on all fine vertices.

\end{proof}

\begin{theorem}[Diagnostic invariance]\label{thm:3.17}

If both sides are $L$-adequate and the comparison of Definition~\ref{def:3.13}
satisfies Condition C, then the actual comparison map $T_{\mathrm{diag}}$ of
\eqref{eq:3.14} is an isomorphism.
In particular, for every coarse diagnostic class $\alpha$,
\begin{equation*}
\alpha=0\quad\Longleftrightarrow\quad T_{\mathrm{diag}}(\alpha)=0
\tag{3.15}\label{eq:3.15}
\end{equation*}
\end{theorem}

\begin{proof}

By Propositions~\ref{prop:3.12} and~\ref{prop:3.14} it suffices to fix one label and
prove injectivity and surjectivity.
C0 ensures that the support of each coarse chart is covered by the images of the
supports of the fine charts in its fiber, and is used throughout as the premise for
comparing the two subnerves of the fixed label over the same support.

\textbf{Injectivity.}
Suppose that the image of a coarse cocycle $z$ satisfies $T^1z=d_f^0b$.
If an edge within a fiber is contracted, then $d_f^0b=0$ on it.
If it is not contracted, it maps to a coarse self-loop, and by C6 it is a self-loop
on the fine side as well, so the values of $b$ at its two endpoints again agree.
By C1, $b$ is constant on each fiber and defines values $\bar b$ at the coarse
vertices.
Taking a lift of each coarse edge by C2 gives
$z(e)=\bar b(t(e))-\bar b(s(e))$.
Hence $z=d_c^0\bar b$, and the original class is zero.

\textbf{Surjectivity.}
Take a fine cocycle $w$. By Lemma~\ref{lem:3.16} the cochain
$w'=w-d_f^0p$, representing the same class, can be made zero within fibers.
By C2 and C5, each coarse edge has exactly one lift with the same label.
Let $z(e)$ be the value of $w'$ there.
On mapped fine edges $T^1z=w'$, and contracted edges lie within fibers, so both sides
are zero there.
Lifting each coarse face by C4, the uniqueness of the three edges shows that
$d_c^1z$ equals $d_f^1w'=0$ on the lifted face.
Hence $z$ is a cocycle and $[w]=T_{\mathrm{diag}}[z]$.

\end{proof}

Condition C is the condition for correcting locally removable differences within
fibers and reading the remaining differences back to the coarse side uniquely.
In the next section we also give a decision that computes the comparison itself,
without using Condition C.

\begin{example}[Comparison from three charts to four]\label{ex:3.18}

Use the two resolutions of Example~\ref{ex:3.5}, with every chart support equal to
the whole set of resolution values.
The nerves are taken from the three-chart and four-chart covers of
Lemma~\ref{lem:2.38}; the two covers and the merging of charts were shown in
Figure~\ref{fig:2.1}.
Map $a_0,a_1$ to $a$, and $b,c$ to the charts of the same name.
Contract the edge $e_{01}$, and map $ab,bc,ac$ to the edges of the same name.

For each label, the fiber over $a$ is a tree joining two vertices by a single edge,
and the fibers over $b,c$ are single points.
The supports agree, the lifts of the three coarse edges are unique, and faces and
coarse self-loops are absent on both sides.
Condition C therefore holds.
Each of the two nerves has one independent cycle and there are two labels 0 and 1,
so both diagnostic groups are $\mathbb{Q}^2$.
The comparison sends the values $(z_{ab},z_{bc},z_{ac})$ on the three edges to
$(0,z_{ab},z_{bc},z_{ac})$.
The period of a fine 1-cochain $w$ is
$w_{e_{01}}+w_{ab}+w_{bc}-w_{ac}$.
On the image of the comparison, the value $w_{e_{01}}$ on the internal edge is zero
and the period agrees with the coarse period $z_{ab}+z_{bc}-z_{ac}$, so nonzero
classes are preserved as well.

\end{example}

\begin{example}[Preserving Laws versus preserving diagnoses]\label{ex:3.19}

Let the sources be $\{0,1,2\}$, the coarse resolution
$q_c(0)=q_c(1)=0,\ q_c(2)=1$, and the fine resolution the identity.
The Law $v_0=q_c$ is readable on both sides.
Take every chart support to be the whole set and consider the following two comparisons.

\begin{xltabular}{\linewidth}{@{}LLL@{}}
\toprule
Coarse nerve & Fine nerve and comparison & Comparison of $H^1$ for one Law value \\
\midrule
\endhead
One vertex, one self-loop, no face & Send a single edge joining two vertices to the self-loop & $\mathbb{Q}\to0$; the nonzero coarse class dies \\
\addlinespace[3pt]
Two vertices, one edge, no face & Send two parallel edges between the two vertices to the single coarse edge & $0\to\mathbb{Q}$; the nonzero fine class is not captured \\
\bottomrule
\end{xltabular}

In the first row, any value on the fine edge is a difference of vertex values, and C6
fails. In the second row, a coboundary takes the same value on the two edges, so
their difference detects the nonzero class, and C5 fails.
In both rows, adequacy with respect to $L=\{0\}$ does hold.

Now add a Law reading $v_1(1)=1,\ v_1(0)=v_1(2)=0$; then $q_c$ can no longer read
this Law.
For an arbitrary resolution, one can still build the diagnosis from the subfamily
$L(q)=\{\ell\in L\mid v_\ell\text{ descends through }q\text{}\}$
of Laws that descend. In this example $L(q_c)=\{0\}$ and
$L(q_f)=\{0,1\}$.
The fine diagnosis of the first row is zero even with the added Law, while in the
second row the nonzero class from the original Law remains.
With a subfamily the computation is defined, but one must check both the retained
Laws and the comparison conditions on the covers in order to judge the correspondence
with the original diagnosis.

\end{example}

\section{Deciding invariance for all Laws}\label{sec:3.6}

Having verified invariance for one Law family, we may ask whether the same comparison
geometry works for every Law family for which both resolutions are adequate.
Instead of enumerating all Law families, it suffices to examine subsets of the values
of the coarse resolution.

\begin{definition}[Nerves selected by subsets of values]\label{def:3.20}

Fix the pair of resolutions and the comparison of Definition~\ref{def:3.13}.
For $A\subseteq Q_c$, retain the coarse cells whose supports meet $A$ and the fine
cells whose supports meet $\pi^{-1}(A)$.
Write these as $N_{c,A}$ and $N_{f,\pi^{-1}(A)}$.
If a face remains, so do its three edges; if an edge remains, so do its two
endpoints.
The image of a mapped fine cell also remains, so the same rule as \eqref{eq:3.13}
yields
\begin{equation*}
T_A:H^1(N_{c,A};\mathbb{Q})
\longrightarrow H^1(N_{f,\pi^{-1}(A)};\mathbb{Q})
\tag{3.16}\label{eq:3.16}
\end{equation*}
Define the pair of dimensions of its kernel and cokernel as
\begin{equation*}
J_A=\left(\dim_{\mathbb{Q}}\ker T_A,\
\dim_{\mathbb{Q}}\bigl(H^1(N_{f,\pi^{-1}(A)};\mathbb{Q})/\mathrm{im}\,T_A\bigr)\right)
\tag{3.17}\label{eq:3.17}
\end{equation*}
The first component is the dimension of the coarse classes killed by the comparison,
and the second is the dimension of the fine classes not captured by the image from
the coarse side.

\end{definition}

\begin{definition}[Uniform invariance]\label{def:3.21}

We say that the comparison is uniformly invariant if, for every finite Law family on
the same sources, the map \eqref{eq:3.14} is an isomorphism whenever both resolutions
are adequate.
The quantification here ranges over Law families including the value sets and the
evaluation maps of the individual Laws.
The resolutions, the nerves, the supports, and the cell comparison maps are fixed.

\end{definition}

\begin{theorem}[Reduction of uniform invariance to finite subsets]\label{thm:3.22}

The following are equivalent.

\begin{enumerate}
\item The comparison is uniformly invariant.
\item For every nonempty $A\subseteq Q_c$, the map $T_A$ is an isomorphism.
\item For every nonempty $A\subseteq Q_c$, we have $J_A=(0,0)$.
\end{enumerate}

\end{theorem}

\begin{proof}

We first show that 2 implies 1. Take a Law family for which both resolutions are
adequate and one label $\lambda=(\ell,v)$, and set
\begin{equation*}
A_\lambda=\{u\in Q_c\mid\bar v_{\ell,c}(u)=v\}
\tag{3.18}\label{eq:3.18}
\end{equation*}
This set is nonempty because the label is realized on the sources.
A coarse cell carries $\lambda$ exactly when its support meets $A_\lambda$.
On the fine side, by Lemma~\ref{lem:3.2}, the corresponding condition is that the
support meets $\pi^{-1}(A_\lambda)$.
The complex and the comparison map for each label therefore coincide with those of
\eqref{eq:3.16}.
By 2 and Proposition~\ref{prop:3.12}, the total comparison is an isomorphism as well.

Conversely, assume 1. For a nonempty $A\subseteq Q_c$, define the evaluation of a
single Law by setting $v_A(s)$ equal to 1 when $q_c(s)\in A$ and to 0 otherwise.
This Law is readable on both resolutions, and the comparison of the subnerves for
its value 1 is exactly $T_A$.
The total comparison is the direct sum over Law values, so its bijectivity implies
the bijectivity of each summand.

The equivalence of 2 and 3 is finite-dimensional linear algebra.
Zero kernel dimension means injectivity and zero cokernel dimension means
surjectivity, and conversely.

\end{proof}

\begin{proposition}[Decision from a finite presentation]\label{prop:3.23}

If finite enumerations of the sources, of the values of both resolutions, and of the
cells, the value tables of both readings, membership in the supports, and the
incidence relations and the cell comparison are given in computable form, then
uniform invariance is decidable.

\end{proposition}

\begin{proof}

For each $u\in Q_f$, find a source representative by finite search and compute
$\pi(u)=q_c(s)$.
Lemma~\ref{lem:3.2} guarantees independence of the representative.
For each nonempty subset $A\subseteq Q_c$, enumerate the remaining cells and build
the rational matrices $d_c^0,d_c^1,d_f^0,d_f^1,T^1$ from \eqref{eq:3.9} and \eqref{eq:3.13}.

By rational elimination, compute bases of $Z_c=\ker d_c^1$, $Z_f=\ker d_f^1$,
$B_c=\mathrm{im}\,d_c^0$, and $B_f=\mathrm{im}\,d_f^0$.
Extending the basis of each $B$ to a basis of the corresponding $Z$ yields a basis
of the quotient.
Expressing $T^1$ in these bases and writing $r_A$ for the rank of the induced map,
we have
\begin{equation*}
J_A=(\dim Z_c-\dim B_c-r_A,\
\dim Z_f-\dim B_f-r_A)
\tag{3.19}\label{eq:3.19}
\end{equation*}
Check whether both components are zero for the finitely many $A$.
The generated matrices are exactly the differentials and the comparison defined above, so
Theorem~\ref{thm:3.22} gives the correctness of the decision.

\end{proof}

This decision does not merely compare the dimensions of the groups on the two sides.
A map between groups of the same dimension can still be the zero map, so we compute
the rank $r_A$ of the actual comparison.

\begin{corollary}[Condition C for all subsets]\label{cor:3.24}

Suppose that C0, C5, and C6 hold globally, and that for every nonempty
$A\subseteq Q_c$ the clauses C1--C4 hold on the two selected subnerves.
Then the comparison is uniformly invariant.
These geometric conditions themselves are decidable from the same finite
presentation as in Proposition~\ref{prop:3.23}.

\end{corollary}

\begin{proof}

For an arbitrary Law family for which both resolutions are adequate, take the sets
$A_\lambda$ of \eqref{eq:3.18}.
The identification used in the proof of Theorem~\ref{thm:3.22} preserves vertices,
edges, faces, their incidences, and the partial maps, so the paths of C1, the lifts
of C2 and C4, and the chains and internal faces of C3 transfer directly to the
subnerves of the labels.
C0, C5, and C6 already hold globally. Uniform invariance follows from
Theorem~\ref{thm:3.17}.

For the finite decision, use finite matching for C0, C2, C4, C5, and C6, and graph
reachability for C1.
For C3, compute bases of $\ker\partial_1$ and $\mathrm{im}\,\partial_2$ for each
fiber and check whether the former is contained in the latter.
Since $\partial_1\partial_2=0$, this is also a test of equality of the two spaces.

\end{proof}

\begin{example}[Condition C is sufficient, and its clauses are not necessary]\label{ex:3.25}

The following finite examples show that uniform invariance can hold while individual
clauses fail.
The resolutions are the $q_c,q_f$ of Example~\ref{ex:3.5}, and every chart support is
the whole set.
To both sides of each row we add, as an independent component, a component with one
vertex, one self-loop $h$, and no face, on which the comparison is the identity.
This common component has nonzero $H^1\cong\mathbb{Q}$.
Adding it makes each comparison an example that also carries a nonzero class.
The table records only the parts added beyond it.

\begin{xltabular}{\linewidth}{@{}LL@{}}
\toprule
Failing clause & Added parts on the coarse and fine sides, and the comparison \\
\midrule
\endhead
C0 & Add one isolated vertex on the coarse side, with no corresponding vertex on the fine side \\
\addlinespace[3pt]
C1 & Map two isolated fine vertices to one isolated coarse vertex \\
\addlinespace[3pt]
C2 & Take a triangle with a face on the coarse side and the path formed by two of its edges on the fine side, compared by the inclusion \\
\addlinespace[3pt]
C3 & Add to both sides a component with one vertex, one self-loop, and no face, compared by the identity \\
\addlinespace[3pt]
C4 & On the coarse triangle place two faces with the same boundary; on the fine side place the same three edges and one face \\
\addlinespace[3pt]
C5 & Take a triangle with a face on the coarse side; on the fine side split one edge into two parallel edges and place one face using each. Send the parallel edges to the same edge \\
\addlinespace[3pt]
C6 & On the coarse side place one self-loop $e$ and a face with boundary $e-e+e=e$. The fine side is a single edge joining two vertices, sent to $e$ \\
\bottomrule
\end{xltabular}

The comparisons in the table send the remaining vertices and edges to those of the
same name.
In the C4 row the fine face is sent to one of the two coarse faces, and in the C5
row the two fine faces are sent to the coarse face.
All satisfy Definition~\ref{def:3.13}.

On the added part of the C3 row the comparison is the identity.
In every other row, $H^1$ of the added part is zero on both sides.
The C0 and C1 rows have no edges, and in the path and the faced triangle of the C2
row every cocycle is a difference of vertex values.
The two faces of the C4 row impose the same equation, so the kernel is the same as
with a single face.
In the C5 row, subtracting the equations of the two faces makes the values on the
parallel edges equal, and the remaining triangle equation then writes the cocycle as
a difference of vertex values.
In the C6 row the coarse face equation imposes $z_e=0$, and the fine side is a path,
hence zero as well.

Since all supports are the whole set, the cells of the two nerves are unchanged for any
choice of nonempty $A$.
Each $T_A$ is the identity on the common nonzero component and the isomorphism above
on the rest.
By Theorem~\ref{thm:3.22} each row is uniformly invariant, which shows that the
sufficient condition of Corollary~\ref{cor:3.24} is strictly stronger than uniform
invariance.
The common self-loop without a face cannot be expressed as a boundary of internal
faces, so C3 also fails in every row.
In this way, clauses other than the listed one may fail simultaneously in each row.
Each row is an example showing that the clause listed in the table is not necessary.

\end{example}

\section{An example undecidable from local incidence information alone}\label{sec:3.7}

The finite decision computed the relations among faces on the complex as a whole.
Reading only the kinds and numbers of incidences locally need not allow the same
decision.
We fix which information is read, and compare two examples on which it agrees.

\begin{definition}[Local observation of incidence information]\label{def:3.26}

Let $\mathrm{Obs}_{\mathrm{loc}}$ be the observation given by the following rules.

\begin{enumerate}
\item Record the sources, the values of both resolutions, and both readings.
For each $A\subseteq Q_c$ with $A\ne\varnothing$, form the two subnerves of
Definition~\ref{def:3.20}.
\item Record as the color of each cell its degree, whether it is coarse or fine, and
its support within the chosen subset.
For a fine cell, also record whether the comparison maps it to a cell of the same
degree or contracts it.
\item For each cell, tally the colors of the directly incident cells by kind of
incidence: the start and the end of an edge, and the 0th, 1st, and 2nd edges of a
face.
The count for each item is collapsed to the three values 0, 1, and 2 or more.
The color of the cell itself together with this tally is the local type of that
cell.
\item The numbers of occurrences of the local types on each side are likewise
collapsed to 0, 1, and 2 or more.
The tables obtained for all $A$ form the observation value.
\end{enumerate}

Incidence is read between charts and edges and between edges and faces, and the
proper names of cells are not recorded.
This is an explicit observation rule retaining the kinds of neighbors and their
truncated counts.

\end{definition}

\begin{construction}[Sequences of faces of length 3 and length 6]\label{cons:3.27}

Take the same $S=\{0,1,2\}$ and resolutions as in Example~\ref{ex:3.19}.
Place charts $w_0,w_1$ on both sides, with supports
\begin{equation*}
\begin{array}{c|cc}
&w_0&w_1\\ \hline
\text{coarse}&\{0\}&\{0,1\}\\
\text{fine}&\{0,1\}&\{2\}
\end{array}
\tag{3.20}\label{eq:3.20}
\end{equation*}
Place one self-loop $h$ at $w_0$ and self-loops $e_0,\ldots,e_{n-1}$ at $w_1$.
Let $f_0,\ldots,f_{n-1}$ be the faces over $w_1$, the three edges of $f_j$ being
$(e_j,e_{j+1},e_{j+2})$, with indices read modulo $n$.
The comparison of charts, edges, and faces is the identity onto the cells of the
same name.
Since \eqref{eq:3.20} satisfies \eqref{eq:3.12}, the comparison is defined.
We call the two inputs with $n=3,6$ respectively $P_3,P_6$.

\end{construction}

\begin{proposition}[Same local observation, different uniform invariance]\label{prop:3.28}

We have $\mathrm{Obs}_{\mathrm{loc}}(P_3)=\mathrm{Obs}_{\mathrm{loc}}(P_6)$.
On the other hand, $P_3$ is uniformly invariant and $P_6$ is not.
Hence no predicate of this observation value alone can decide uniform invariance.

\end{proposition}

\begin{proof}

We first compare the observations. The component of $h$ is identical in the two
inputs.
When the component of $w_1$ remains, each $e_j$ has charts of the same color at both
endpoints and is incident to exactly one face in each of the three edge positions of
a face.
Each face is likewise incident to edges of the same color in its three positions.
The number of edges incident to $w_1$, and the numbers of occurrences of the local
types of edges and faces, are all 2 or more for both $n=3,6$.
For each subset, which sides retain this component is determined by \eqref{eq:3.20}
and does not depend on $n$.
All observation tables therefore agree.

We next compute the cohomology. All edges are self-loops, so $d^0=0$.
The cocycle condition over $w_1$ is
\begin{equation*}
z_j-z_{j+1}+z_{j+2}=0
\qquad(j\ \bmod n)
\tag{3.21}\label{eq:3.21}
\end{equation*}
Starting from $z_0=a,z_1=b$, the sequence is
\begin{equation*}
a,\ b,\ b-a,\ -a,\ -b,\ a-b,\ a,\ b,\ldots
\tag{3.22}\label{eq:3.22}
\end{equation*}
Imposing period 3 gives $a=-a,\ b=-b$ and hence $a=b=0$, while with period 6
arbitrary $a,b\in\mathbb{Q}$ are allowed.
The dimension of $H^1$ of this component is therefore 0 for $n=3$ and 2 for $n=6$.

For $A=\{0\}$, both components remain on the coarse side and only the component of
$h$ remains on the fine side.
The map $T_A$ preserves the value on $h$ and discards the component of $w_1$.
For $A=\{1\}$, only the component of $w_1$ remains on both sides, and the comparison
is the identity.
For $A=\{0,1\}$, all components remain on both sides, and the comparison is again
the identity.
Hence
\begin{equation*}
\begin{array}{c|ccc}
&A=\{0\}&A=\{1\}&A=\{0,1\}\\ \hline
P_3:\ J_A&(0,0)&(0,0)&(0,0)\\
P_6:\ J_A&(2,0)&(0,0)&(0,0)
\end{array}
\tag{3.23}\label{eq:3.23}
\end{equation*}
and Theorem~\ref{thm:3.22} gives the conclusion.
If uniform invariance were determined by a predicate of the observation value alone,
these two inputs, whose observations agree, would receive the same truth value.
This contradicts the computation above.

\end{proof}

What produces the difference is the period obtained by following the face relations
around a full turn.
In this example, a relation that cannot be read off from the agreement of local
types is retained by the totality of \eqref{eq:3.21}.

\section{Structural supports and the role of coefficients}\label{sec:3.8}

Chapter~\ref{chap:2} compared the obstructions of the structural and semantic
phases.
Here we directly build the structural support that survives changes of semantic
choices, and study how the choice of coefficients placed on it affects the
diagnosis.
To this end we generate a nerve from all tuples of Atoms holding at a common source,
and also separate the coefficients by source.
In this construction the correction from a reference Atom can be made explicit, and
we can trace the mechanism by which all diagnostic classes become zero.

\begin{definition}[Supports invariant under semantic changes]\label{def:3.29}

For an extraction doctrine of Chapter~\ref{chap:1}, fix the vocabulary, the
resolution, the semantics of sources, and the normalization, and vary only the
semantic reading and its permission predicates. Specify these options by a set
$\Theta$.
Designate a reference choice $\theta_0\in\Theta$, and let $E_\theta(s,a)$ mean that
``Atom $a$ is extracted from source $s$''.
Define the structural support by
\begin{equation*}
E_{\mathrm{str}}(s,a)
\quad\Longleftrightarrow\quad
E_{\theta_0}(s,a)\ \land\
\forall\theta\in\Theta,\
\bigl(E_\theta(s,a)\Longleftrightarrow E_{\theta_0}(s,a)\bigr)
\tag{3.24}\label{eq:3.24}
\end{equation*}
The nerve built from a support $E(s,a)$ has as cells the ordered tuples of Atoms
holding at a common source.
In degree 0 take $(a)$, in degree 1 take $(a,b)$, and in degree 2 take $(a,b,c)$,
the condition being the existence of a source satisfying all entries simultaneously.
Repeated Atoms are allowed, and the incidences of edges and faces are given by the
projections of the tuples.

\end{definition}

\begin{proposition}[Invariance of the structural nerve]\label{prop:3.30}

The nerve obtained from $E_{\mathrm{str}}$ has the same cells and incidence
relations when the reference semantic reading is replaced by another member of the
same family.

\end{proposition}

\begin{proof}

\eqref{eq:3.24} is equivalent to $\forall\theta\in\Theta,\ E_\theta(s,a)$, and hence
does not depend on the reference member. Since the support is the same, the
vertices, edges, and faces defined by the existence of a common source, and the
projections of the tuples, all coincide.

\end{proof}

\begin{example}[The all-Atom nerve can change]\label{ex:3.31}

Consider the case of one source and two Atoms $a,b$.
Suppose that one semantic reading extracts only $a$ and the other extracts $a,b$.
Only $a$ remains in the structural support, while the all-Atom nerve acquires the
vertex $b$ under the latter.
Invariance of the structural part is distinguished from the conclusion that the
whole nerves are equal.

\end{example}

\begin{proposition}[Vanishing with coefficients separated by source]\label{prop:3.32}

Suppose that $S$ and the set of Atoms are finite.
On each cell of the ordered nerve of Definition~\ref{def:3.29}, place one rational
coordinate for each source realizing that cell simultaneously.
The differentials leave the source unchanged and are given by
\begin{equation*}
\begin{aligned}
(d^0c)(a,b;s)&=c(b;s)-c(a;s),\\
(d^1z)(a,b,c;s)&=z(b,c;s)-z(a,c;s)+z(a,b;s)
\end{aligned}
\tag{3.25}\label{eq:3.25}
\end{equation*}
The $H^1$ of this complex is zero.

\end{proposition}

\begin{proof}

For each source $s$ with nonempty support, choose one Atom satisfying $E(s,a_s)$.
For a cocycle $z$, set $c(a;s)=z(a_s,a;s)$.
If the edge $(a,b;s)$ exists, then so does the face $(a_s,a,b;s)$, and its cocycle
condition gives $z(a,b;s)=z(a_s,b;s)-z(a_s,a;s)$.
That is, $z=d^0c$. A source with empty support has no coordinates.

\end{proof}

In this construction, the coordinates sharing a source can each be corrected from
the reference Atom.
To diagnose a nonzero gluing obstruction, therefore, one must examine the relation
between these coefficients separated by source and the actual obstruction
coefficients.
In the next section we give this comparison, using the integer coefficients built
from the generator relations of Chapter~\ref{chap:2}.

\section{Mapping integer obstruction coefficients to the Law-value diagnosis}\label{sec:3.9}

The obstructions of Chapter~\ref{chap:2} arose from differences obtained by actually
restricting local states and comparing them along transitions.
In this section we read the Law values of the generators composing those
differences and transfer them to the diagnostic coordinates.
We construct in turn the map of coefficients, the map of cochains, and the map of
obstruction classes.

\begin{construction}[Coefficient comparison]\label{cons:3.33}

Take the generators $G=L\times S$ of Definition~\ref{def:2.35}, the label-preserving
primitive relation $R$, the set $\mathfrak{B}$ of relation components, and the
integer coefficients $M_R$.
Write $\lambda_B\in\Lambda$ for the common label of a component $B\in\mathfrak{B}$.
Send the class of a generator to the rational delta function of its label:
\begin{equation*}
\varepsilon_R:M_R\longrightarrow \mathbb{Q}^\Lambda,\qquad
\varepsilon_R([e_g])(\lambda)
=\begin{cases}
1&\lambda=\lambda_g,\\
0&\lambda\ne\lambda_g.
\end{cases}
\tag{3.26}\label{eq:3.26}
\end{equation*}
If $gRh$ then the labels are equal and $e_g-e_h$ maps to zero, so this is well
defined as an additive homomorphism on the quotient.
Under the presentation $M_R\cong\mathbb{Z}^{(\mathfrak{B})}$ of
Lemma~\ref{lem:2.36},
\begin{equation*}
\varepsilon_R(m)(\lambda)
=\sum_{\lambda_B=\lambda}m_B
\quad\text{(integers read as rationals via the inclusion)}
\tag{3.27}\label{eq:3.27}
\end{equation*}
The coefficients of the relation components carrying the same Law value are added
into a single diagnostic coordinate.

\end{construction}

\begin{lemma}[Injectivity of the coefficient comparison]\label{lem:3.34}

Suppose that any two generators with the same label are joined by a finite chain of
primitive relations.
Then $B\mapsto\lambda_B$ gives $\mathfrak{B}\cong\Lambda$, and
$\varepsilon_R$ is injective.

\end{lemma}

\begin{proof}

Every label admits a generator, so the map is surjective.
If the labels of two components are equal, the hypothesis joins their representative
generators by relations, and the components are equal as well.
Hence \eqref{eq:3.27} becomes $\varepsilon_R(m)(\lambda_B)=m_B$.
Since the inclusion of the integers into the rationals is injective, all coordinates
can be read back.

\end{proof}

\begin{example}[Different components sent to the same Law value]\label{ex:3.35}

Take two sources, one constant Law, and the empty primitive relation.
Then $M_R=\mathbb{Z}^2$, the set $\Lambda$ is a singleton, and
$\varepsilon_R(m,n)=m+n$.
The nonzero element $(1,-1)$ maps to zero.
Making ``reading the same Law value'' coincide with ``lying in the same component of
the primitive relation'' is the role of Lemma~\ref{lem:3.34}.

\end{example}

\begin{construction}[Comparison with the actual \v{C}ech complex]\label{cons:3.36}

Take a finite cover $\mathcal{U}_q$ consisting of nonempty connected charts, in which the
nonempty double intersections are connected and the intersections of three distinct
charts are empty.
The nerve has as edges the nonempty intersections of distinct pairs of charts,
oriented by the order of the charts.
On this cover place the locally constant sheaf $F_R$ of $M_R$-values from
Construction~\ref{cons:2.39}.

By connectedness, each section is determined by a single $M_R$-value, and the
restrictions preserve that value.
The normalized \v{C}ech complex used in Chapter~\ref{chap:2} can therefore be
presented as
\begin{equation*}
C_{\mathrm{ob}}^0(q)=M_R^{N_{q,0}},\qquad
C_{\mathrm{ob}}^1(q)=M_R^{N_{q,1}},\qquad
C_{\mathrm{ob}}^2(q)=0
\tag{3.28}\label{eq:3.28}
\end{equation*}
Here $d_{\mathrm{ob}}^0p$ is the difference, after the actual restrictions, between
the value on the end side and the value on the start side of an edge.
On the same nerve, specify the chart supports of an $L$-adequate $q$ and build the
diagnostic complex of Construction~\ref{cons:3.10}.
Define the comparisons in degrees 0 and 1 by
\begin{equation*}
(\phi_q^n c)_{\sigma,\lambda}
=\varepsilon_R(c_\sigma)(\lambda)
\quad(n=0,1),\qquad
\phi_q^2:0\longrightarrow0
\tag{3.29}\label{eq:3.29}
\end{equation*}
On the diagnostic side, only the labels occurring on the cell are read.
The additivity of $\varepsilon_R$ and the presentation of the restrictions give
\begin{equation*}
\phi_q^1d_{\mathrm{ob}}^0=d_{\mathrm{diag}}^0\phi_q^0,\qquad
\phi_q^2d_{\mathrm{ob}}^1=d_{\mathrm{diag}}^1\phi_q^1
\tag{3.30}\label{eq:3.30}
\end{equation*}
the latter because degree 2 is zero on both sides.
We therefore obtain an additive homomorphism
\begin{equation*}
\Phi_q:\check H^1(\mathcal{U}_q,F_R)\longrightarrow
H^1_{\mathrm{diag}}(q,N_q;L),\qquad
[z]\longmapsto[\phi_q^1z]
\tag{3.31}\label{eq:3.31}
\end{equation*}
The domain is an abelian group with integer coefficients, and the structure of this
comparison is that of an additive homomorphism.

\end{construction}

\begin{theorem}[Obstruction and diagnostic classes from the same local data]\label{thm:3.37}

For the input of Construction~\ref{cons:3.36}, take any local states and affine
transitions $x=(\xi,p)$.
Define the obstruction cochain of Proposition~\ref{prop:2.41} and the diagnostic
cochain generated from Law values by
\begin{equation*}
\begin{aligned}
o_q(x)&=\xi+d_{\mathrm{ob}}^0p,\\
a_q(x)&=\phi_q^1\xi+d_{\mathrm{diag}}^0\phi_q^0p
\end{aligned}
\tag{3.32}\label{eq:3.32}
\end{equation*}
Then
\begin{equation*}
\phi_q^1o_q(x)=a_q(x),\qquad
\Phi_q([o_q(x)])=[a_q(x)]
\tag{3.33}\label{eq:3.33}
\end{equation*}
In particular, $[o_q(x)]=0$ implies $[a_q(x)]=0$.

\end{theorem}

\begin{proof}

The value of the diagnosis at an edge coordinate $(e,\lambda)$ is the sum of the
Law-value coordinate $\varepsilon_R(\xi_e)(\lambda)$ of the transition and the
difference of the Law-value coordinates of the states on the two sides.
By \eqref{eq:3.30} and additivity, this equals $\varepsilon_R(o_q(x)_e)(\lambda)$.
Degree 2 is zero on both sides, so both cochains are cocycles, and passing to the
quotients gives \eqref{eq:3.33}.

\end{proof}

In this comparison, we first prepare the integer transitions and states, and then
read their Law values to build the diagnosis.
Thus \eqref{eq:3.33} does more than compare two groups abstractly: it determines the
destination of the concrete obstruction classes received from
Chapter~\ref{chap:2}.

\section{From vanishing diagnoses to integer corrections}\label{sec:3.10}

In the reverse direction, a correction seen in the diagnosis must be brought back to
a correction in the original integer coefficients.
For this we use the facts that the relation components can be distinguished by Law
values, and that the required Law values can be read on every chart and every edge.

\begin{definition}[Two conditions reflecting vanishing]\label{def:3.38}

Impose the following conditions on the input of Construction~\ref{cons:3.36}.

\begin{itemize}
\item \textbf{$R_q$: identification of relation components.} Any two generators with
the same label are joined by a finite chain of primitive relations.
\item \textbf{Existence of common representatives.} For each
$\lambda=(\ell,v)\in\Lambda$ there is a single $u_\lambda\in Q$ that belongs to
every chart support and satisfies $\bar v_\ell(u_\lambda)=v$.
\end{itemize}

By the second condition, that representative belongs to the support of every edge as
well.
Each label therefore has one diagnostic coordinate on every chart and every edge.
Even if the representatives are chosen differently, the coordinates of
Construction~\ref{cons:3.10} are treated as the same label.

\end{definition}

\begin{lemma}[Integer corrections from rational ones]\label{lem:3.39}

Let an integer-valued 1-cochain $z$ and a rational-valued 0-cochain $b$ on a
directed graph satisfy $b_{t(e)}-b_{s(e)}=z_e$.
Setting $n_\alpha=\lfloor b_\alpha\rfloor$ gives
$n_{t(e)}-n_{s(e)}=z_e$.

\end{lemma}

\begin{proof}

From $b_{t(e)}=b_{s(e)}+z_e$ and $z_e\in\mathbb{Z}$ it follows that
$\lfloor b_{t(e)}\rfloor=\lfloor b_{s(e)}\rfloor+z_e$.

\end{proof}

For example, the difference of $1/2,3/2$ is 1, and the integer parts $0,1$ have the
same difference.
The integer parts are used to construct a correction preserving these specific edge
differences.

\begin{theorem}[Reflection of vanishing for specified obstruction classes]\label{thm:3.40}

Under the two conditions of Definition~\ref{def:3.38}, for arbitrary local data
$x=(\xi,p)$ we have
\begin{equation*}
[a_q(x)]=0
\quad\Longleftrightarrow\quad
[o_q(x)]=0
\tag{3.34}\label{eq:3.34}
\end{equation*}
\end{theorem}

\begin{proof}

Right to left is Theorem~\ref{thm:3.37}.
Assuming the left side, there is a diagnostic 0-cochain $b$ with
$a_q(x)=d_{\mathrm{diag}}^0b$.
Write $z_{e,B}$ for the integer coefficient of the relation component $B$ in
$o_q(x)_e$.
By the common representatives the coordinate $\lambda_B$ exists on the edges, and by
$R_q$ and Lemma~\ref{lem:3.34} we get
\begin{equation*}
\begin{aligned}
b_{t(e),\lambda_B}-b_{s(e),\lambda_B}
&=a_q(x)_{e,\lambda_B}\\
&=\varepsilon_R(o_q(x)_e)(\lambda_B)
=z_{e,B}
\end{aligned}
\tag{3.35}\label{eq:3.35}
\end{equation*}
Applying Lemma~\ref{lem:3.39} to each component and setting
\begin{equation*}
n_\alpha=\sum_{B\in\mathfrak{B}}
\lfloor b_{\alpha,\lambda_B}\rfloor e_B\in M_R
\tag{3.36}\label{eq:3.36}
\end{equation*}
we obtain $d_{\mathrm{ob}}^0n=o_q(x)$.
Since $\mathfrak{B}$ is finite, this sum defines an element of $M_R$.
On the connected charts of Construction~\ref{cons:3.36}, $n_\alpha$ gives an actual
locally constant section.
The original obstruction class is therefore zero as well.

\end{proof}

This proof produces the correction itself.
Changing the original states to $p-n$ gives
$\xi+d_{\mathrm{ob}}^0(p-n)=0$, and the local states become consistent with respect
to the transitions.
In the sheaf of states of Proposition~\ref{prop:2.41}, the corrected family glues to
a global state.

\begin{example}[Without the reflection conditions the diagnosis loses differences]\label{ex:3.41}

On the three-chart cover of Lemma~\ref{lem:2.38}, place the locally constant
coefficients $M_R=\mathbb{Z}^2$ of Example~\ref{ex:3.35} and chart supports reading
all values of the resolution.
Taking $p=0$, $\xi_{ab}=(1,-1)$, and $\xi_{bc}=\xi_{ac}=0$, the period of the
obstruction is nonzero while the diagnostic cochain is zero. This is an example
lacking $R_q$.

The role of the common representatives can be checked in the same fashion.
This time use two Law values 0 and 1, with an $M_R\cong\mathbb{Z}^2$ satisfying
$R_q$.
Taking each chart support to be a set on which only the value 0 occurs, the
coordinate of value 1 does not appear in the diagnosis.
Placing the basis element $u$ of value 1 in $\xi_{ab}$ and zero elsewhere, the
period of the obstruction is $u\ne0$ while the diagnosis is zero.
Even when the coefficient comparison is injective, the reflection of vanishing can
break if the coordinates reading a component are locally missing.

\end{example}

\section{Refinement of covers and invariance of the obstruction verdict}\label{sec:3.11}

Finally, we tie the preceding comparisons together over the same Atom input.
We use the eight-point space and generator relations of Chapter~\ref{chap:2} and the
two resolutions of Example~\ref{ex:3.5}.
The integer obstruction, the rational diagnosis, and the refinement from three
charts to four then appear as one commutative comparison.

\begin{construction}[Coarse--fine comparison from the same input]\label{cons:3.42}

Take the sources $S=\{0,1\}^2$, the Law $v_\ell(v,h)=v$, and the primitive relation
consisting of the two pairs joining $(v,0)$ and $(v,1)$.
The coefficients are $M_R\cong\mathbb{Z}e_{B_0}\oplus\mathbb{Z}e_{B_1}$.
Use the point Atoms and generator Atoms of Construction~\ref{cons:2.37}; each
context reads the points of the corresponding open set and all generators.

On the coarse side use $q_c(v,h)=v$ and the three-chart cover
$\mathcal{U}_c=(U_a,U_b,U_c)$; on the fine side use $q_f=\mathrm{id}_S$ and the
four-chart cover $\mathcal{U}_f=(U_{a_0},U_{a_1},U_b,U_c)$.
Every chart reads all generators, so the diagnostic chart supports are all of
$Q_c,Q_f$, respectively.
Below we abbreviate the obstruction group as
$H^1_{\mathrm{ob}}(q)=\check H^1(\mathcal{U}_q,F_R)$ and the diagnostic group as
$H^1_{\mathrm{diag}}(q)=H^1_{\mathrm{diag}}(q,N_q;L)$.

The inclusions $U_{a_0},U_{a_1}\subseteq U_a$, together with the identity inclusions
of the remaining charts, give an actual refinement of covers.
In the coordinates of \eqref{eq:3.28}, the maps by restriction of sections are
\begin{equation*}
\begin{aligned}
T_{\mathrm{ob}}^0(p_a,p_b,p_c)&=(p_a,p_a,p_b,p_c),\\
T_{\mathrm{ob}}^1(z_{ab},z_{bc},z_{ac})&=(0,z_{ab},z_{bc},z_{ac}),\\
T_{\mathrm{ob}}^2&:0\longrightarrow0
\end{aligned}
\tag{3.37}\label{eq:3.37}
\end{equation*}
with the fine edges ordered $(e_{01},ab,bc,ac)$.
The value on the internal edge $e_{01}$ is zero, being a self-comparison of the same
coarse chart.
On the remaining edges, restriction to the actual overlaps preserves the same
$M_R$-value.

The fine differential is
\begin{equation*}
d_f^0p=
(p_{a_1}-p_{a_0},\
p_b-p_{a_1},\
p_c-p_b,\
p_c-p_{a_0})
\tag{3.38}\label{eq:3.38}
\end{equation*}
so $T_{\mathrm{ob}}^1d_c^0=d_f^0T_{\mathrm{ob}}^0$ can be checked directly.
This yields $T_{\mathrm{ob}}:H^1_{\mathrm{ob}}(q_c)\to H^1_{\mathrm{ob}}(q_f)$.
On the diagnostic side we use the map $T_{\mathrm{diag}}$ generated from the same
cell comparison as in Example~\ref{ex:3.18}.

\end{construction}

\begin{theorem}[Commutativity of obstruction, diagnosis, and refinement]\label{thm:3.43}

In Construction~\ref{cons:3.42} the following square commutes.
\begin{equation*}
\begin{array}{ccc}
H^1_{\mathrm{ob}}(q_c)&\xrightarrow{\ \Phi_{q_c}\ }&
H^1_{\mathrm{diag}}(q_c)\\
\big\downarrow T_{\mathrm{ob}}&&\big\downarrow T_{\mathrm{diag}}\\
H^1_{\mathrm{ob}}(q_f)&\xrightarrow{\ \Phi_{q_f}\ }&
H^1_{\mathrm{diag}}(q_f)
\end{array}
\tag{3.39}\label{eq:3.39}
\end{equation*}
For any coarse local data $x=(\xi,p)$, setting
$x_f=(T_{\mathrm{ob}}^1\xi,T_{\mathrm{ob}}^0p)$ gives
\begin{equation*}
\begin{aligned}
T_{\mathrm{ob}}[o_{q_c}(x)]&=[o_{q_f}(x_f)],\\
T_{\mathrm{diag}}[a_{q_c}(x)]&=[a_{q_f}(x_f)]
\end{aligned}
\tag{3.40}\label{eq:3.40}
\end{equation*}
\end{theorem}

\begin{proof}

The coefficient comparisons on the two sides send the same generator $[e_g]$ to the
delta function of the same Law value.
On mapped edges and charts, applying \eqref{eq:3.37} first and then reading Law
values yields the same coordinates as reading Law values first and then pulling back
by \eqref{eq:3.13}.
On the internal edge both are zero, and degree 2 is zero on both sides as well.
The comparisons therefore commute in each degree, and passing to classes gives
\eqref{eq:3.39}.

For the local data, commutation with the differentials gives
\begin{equation*}
o_{q_f}(x_f)
=T_{\mathrm{ob}}^1\xi+d_f^0T_{\mathrm{ob}}^0p
=T_{\mathrm{ob}}^1(\xi+d_c^0p)
\tag{3.41}\label{eq:3.41}
\end{equation*}
This gives the first equation of \eqref{eq:3.40}.
The second follows from the first, Theorem~\ref{thm:3.37}, and \eqref{eq:3.39}.

\end{proof}

\begin{theorem}[Invariance of the obstruction verdict along a change of reading]\label{thm:3.44}

For any local data of Construction~\ref{cons:3.42}, the following four statements
are equivalent.
\begin{equation*}
\begin{aligned}
[o_{q_c}(x)]=0
&\quad\Longleftrightarrow\quad[a_{q_c}(x)]=0\\
&\quad\Longleftrightarrow\quad[a_{q_f}(x_f)]=0\\
&\quad\Longleftrightarrow\quad[o_{q_f}(x_f)]=0.
\end{aligned}
\tag{3.42}\label{eq:3.42}
\end{equation*}

\end{theorem}

\begin{proof}

For the coarse and the fine resolution alike, the Law value can be read from the
first component, so both are adequate.
Two generators with the same Law value are joined by one of the declared relations,
so $R_q$ holds.
As a common representative of a value $v$ we may take $v$ on the coarse side and
$(v,0)$ on the fine side; these belong to the supports of all charts.
Theorem~\ref{thm:3.40} therefore applies on both ends.

The diagnostic comparison satisfies Condition C, as verified in
Example~\ref{ex:3.18}, and is an isomorphism by Theorem~\ref{thm:3.17}.
Theorem~\ref{thm:3.43} matches the specified diagnostic classes, so their vanishing
is equivalent as well.
Composing these three equivalences gives \eqref{eq:3.42}.

\end{proof}

In this proof, the merging map $\pi(v,h)=v$ is not injective.
Nevertheless, because the Law values, the relation components, and the comparison of
covers satisfy the conditions above, the obstruction verdict for the same local data
can be shared by the two resolutions.

\begin{example}[A nonzero difference that can be removed, and one that cannot]\label{ex:3.45}

Let $u=e_{B_0}\in M_R$ and fix all chart states to zero.
With the same coefficients, Law, and covers, vary only the transitions as follows.

\begin{xltabular}{\linewidth}{@{}LLLL@{}}
\toprule
Local data & Coarse transitions $(ab,bc,ac)$ & Transitions carried to fine $(e_{01},ab,bc,ac)$ & Obstruction and diagnostic classes \\
\midrule
\endhead
A correctable difference & $(u,-u,0)$ & $(0,u,-u,0)$ & Zero on both sides \\
\addlinespace[3pt]
A difference surviving correction & $(u,0,0)$ & $(0,u,0,0)$ & Nonzero on both sides \\
\bottomrule
\end{xltabular}

The first row is $d^0(0,u,0)$ on the coarse side and
$d^0(0,0,u,0)$ on the fine side.
The difference itself is thus nonzero, but adding the negative of the corresponding
0-cochain to the local states restores consistency.

In the second row the coarse period is $u$, and the fine period is $u$ as well.
The fine period is
$z_{e_{01}}+z_{ab}+z_{bc}-z_{ac}$, and substituting \eqref{eq:3.38} shows that it is
zero on every coboundary.
On the other hand, the coordinate of $\varepsilon_R(u)$ at the value 0 is 1, so the
diagnostic period is nonzero as well.
This gives the nonzero verdicts in the table.

In the presentation by periods, the groups and the comparison of this example are
$H^1_{\mathrm{ob}}\cong\mathbb{Z}^2$, $H^1_{\mathrm{diag}}\cong\mathbb{Q}^2$, and
$\Phi:\mathbb{Z}^2\hookrightarrow\mathbb{Q}^2$.
While keeping the difference between integer and rational coefficients,
\eqref{eq:3.42} decides the same vanishing.

\end{example}

\subsection*{Transporting coefficients and local contexts by isomorphisms}

Another basic comparison is the case in which the chosen local contexts and
coefficients are carried by isomorphisms.
If the covers and their intersection diagrams correspond, and the isomorphisms of
coefficients over each intersection commute with the restrictions, then the
comparison of sections gives isomorphisms of the \v{C}ech complexes in each degree.
The differentials are alternating sums of restrictions, hence commute with the
comparison, and the inverse coefficient isomorphisms yield an inverse cochain map.
The vanishing of the corresponding obstruction classes is therefore equivalent.

When the coefficients are built from a ring and an ideal, a ring isomorphism
matching the ideals that define the coefficients induces an isomorphism of the
quotients as well. Matching also the Laws, the witnesses, and the readouts of the
axes can place the classes built from specified residuals into the same comparison.
The comparison by refinement in this chapter goes beyond the case of direct
degreewise isomorphisms: it gives a way to preserve the obstruction verdict through
Condition C and the integer-correction construction.

\section*{Summary of the Chapter}

In this chapter we connected the retained Laws, the local structure, the
coefficients, and the specified local data, and constructed diagnoses along changes
of resolution.

\begin{itemize}
\item \textbf{Canonical resolution.} The quotient merging sources on which all Law
values agree can be read uniquely through every $L$-adequate resolution. Whether it
can be presented by specified extraction methods is determined by comparison with
the equivalence relations of the extraction results.
\item \textbf{Comparison of diagnoses.} We built coefficients from the distinct Law
values occurring on cells, and generated maps commuting with the differentials from
local comparisons. Under Condition C, the map is an isomorphism on $H^1$.
\item \textbf{Uniform invariance.} That the comparison is an isomorphism for every
Law family for which both resolutions are adequate is equivalent to the kernel and
cokernel dimensions being zero for all nonempty subsets of values, and is decidable
from a finite presentation. Condition C is a sufficient condition, and there are
examples on which a fixed local observation agrees while the verdicts differ.
\item \textbf{The role of coefficients.} Even when the structural support is
invariant, which coefficients read the differences is specified independently. For
the coefficients generating all tuples per source, we showed that $H^1$ is zero, by
an explicit formula for the correction.
\item \textbf{Correspondence of obstruction and diagnosis.} We built a map from the
integer generator relations to the Law-value coefficients, and matched the
obstruction and diagnostic classes of the same local data. Under the conditions of
identification of relation components and of common representatives, rational
corrections can be made integral and vanishing is reflected.
\item \textbf{Concrete invariance.} We constructed the comparison square for the
refinement from three charts to four and, on the same input with the non-injective
comparison map $\pi$, carried zero and nonzero obstruction verdicts to both sides.
\end{itemize}

In the opening example, merging internal values that the Law does not read is a
change of resolution preserving the evaluation.
When the regions on which local states are placed are repartitioned as well, a map
built from the comparison of overlaps also becomes necessary.
The results of this chapter give the conditions under which that map preserves
obstruction classes and vanishing, and delimit the range over which the same verdict
can be shared.

\subsection*{Potential applications to code review}

The results of this chapter provide a starting point for thinking about what to
retain in a summary and at what granularity to check during code review.
Consider, for example, a change to order processing in which the inventory update
and the order persistence are reviewed separately.
Besides the description of each process, we also want to check that values and
assumptions are handed over so that the two handle the same order and the same
processing result.
Can the per-function implementations and a description summarizing the handoffs
between modules capture the same inconsistency?

To treat this question with the constructions of this chapter, fix a version and a
semantics of the target code, and define the sources and the Law evaluations from
states and behavior.
The canonical resolution gives the distinctions that a summary must retain in order
to preserve the chosen evaluations.
Taking the reviewed regions as charts, if their overlaps and the correspondences
between regions can be expressed by the nerves and comparison maps of this chapter,
then Condition C of Theorem~\ref{thm:3.17} lets us confirm the preservation of
diagnostic classes along splits and mergers.

Moreover, when obstruction and diagnosis are built from the same local data and the
constructions and hypotheses of \S\S\ref{sec:3.9}--\ref{sec:3.10} are satisfied, a
vanishing diagnosis on the summary can be brought back to the existence of a
correction that makes the original local states consistent.
What vanishing asserts here is correctability under the chosen Laws and
coefficients.
How to realize that correction as an actual code change is a task to be settled at
application time.

The first thing to verify is the agreement of diagnoses when the same code is read
at different granularities.
Checking concretely the correspondence from code to model and the hypotheses of the
comparison leads to a method for evaluating summaries and partitions for review.
To advance to applications comparing the programs before and after a change, one
adds a construction matching their inputs and the specified obstruction classes.

For the gluing obstructions of Chapter~\ref{chap:2}, this chapter has clarified what
is read, and under which conditions the way of reading can be changed while the
same verdict is obtained. The next chapter treats the construction that transports
objects and morphisms along a change of reading, and studies its universality and
its coherence with units and composition.
